
\subsubsection{2.1 Static Analysis for LLM-Generated Code}

Recent studies integrate static analysis into LLM code generation pipelines as feedback rather than predictors. Blyth et al. (2025) demonstrate an iterative feedback loop where Bandit and Pylint errors are fed back to the LLM for self-correction, reducing security vulnerabilities from over 40\% to 13\% and readability violations from over 80\% to 11\% within 10 iterations on HumanEval/MBPP. However, they measure improvement magnitude rather than predictive correlation.

AutoSafeCoder \citep{Nunez2024AutoSafeCoder} combines multi-agent static analysis with fuzz testing, achieving 13\% reduction in code vulnerabilities. CodeQUEST \citep{Liu2025CodeQUEST} uses GPT-4o to iteratively improve code quality as measured by Pylint Score, Radon Maintainability, and Bandit logs, reporting 52.6\% mean improvement and ''meaningful correlation'' between static analysis metrics and quality. Neither study quantifies this correlation with r-values or regression coefficients.

STALL+ \citep{Liu2024STALL} integrates static analysis at the prompting phase for repository-level code completion, finding that static analysis integration performs best when combined with retrieval-augmented generation. Their focus is on improving generation quality through static-analysis-informed prompts, not on using static analysis metrics as standalone predictors.

This work differs in that it measures correlation directly, enabling static-analysis-based rejection sampling without any LLM feedback loop.

\subsubsection{2.2 Self-Repair and Iterative Refinement}

The self-repair paradigm feeds execution feedback back to LLMs for correction. Olausson et al. (2024) systematically study self-repair effectiveness, finding it improves pass rates but is not universally effective across models. Arimbur (2026) shows self-repair improves HumanEval pass rates by +4.9 to +17.1 percentage points, with most gains occurring in the first 2 iterations.

CodeCoR \citep{Pan2025CodeCoR} achieves 77.13\% average Pass@1 on HumanEval/MBPP using multi-agent pruning and test case generation. INTERVENOR \citep{NEUIR2024} employs Code Teacher and Code Learner agents with compiler feedback.

These approaches assume static analysis feedback is useful for correction. The present work tests whether static analysis signal predicts correctness in the first place. The finding that r=0.873 supports the assumption underlying these iterative approaches.

\subsubsection{2.3 Code Generation Evaluation}

Chen et al. (2021) introduced HumanEval with 164 hand-crafted problems and the pass@k metric. Austin et al. (2021) created MBPP with 974 problems. Liu et al. (2023) extended these with EvalPlus, adding $80\times$ more test cases and revealing that 19-29\% of previously ''correct'' code fails under rigorous testing.

Yetistiren et al. (2023) evaluate Copilot, CodeWhisperer, and ChatGPT on HumanEval across correctness, security, reliability, and maintainability dimensions, using Radon, Bandit, and Pylint as quality metrics. However, they report quality scores and correctness separately without correlation analysis.

\begin{table}[htbp]
\centering
\resizebox{\columnwidth}{!}{%
\begin{tabular}{llll}
\toprule
Work & SA Usage & Correlation Measured & r-Value Reported \\
\midrule
Blyth et al. (2025) & Feedback loop & No & --- \\
CodeQUEST (2025) & Iterative improvement & ''Meaningful'' (qualitative) & --- \\
AutoSafeCoder (2024) & Multi-agent feedback & No & --- \\
Yetistiren et al. (2023) & Quality measurement & No & --- \\
This work & Prediction & Yes & r=0.873 \\
\bottomrule
\end{tabular}
}
\end{table}
